\documentclass[11pt]{article}

\usepackage{amsmath,amssymb}
\usepackage{xurl}
\usepackage[hidelinks]{hyperref}
\setlength{\emergencystretch}{2em}

% Shared commands for notes in docs/paper-gaps/.
% Notation follows the LDT paper (references/ldt-paper/).

\newcommand{\Lean}{\textsc{Lean}}
\newcommand{\leanid}[1]{\textsf{#1}}
\newcommand{\ghissue}[1]{\href{https://github.com/LionSR/MIPStarRE/issues/#1}{GitHub issue \##1}}
\newcommand{\ghpr}[1]{\href{https://github.com/LionSR/MIPStarRE/pull/#1}{GitHub PR \##1}}

% --- Paper-compatible notation ---
\newcommand{\F}{\mathbb{F}}
\newcommand{\N}{\mathbb{N}}
\newcommand{\C}{\mathbb{C}}
\newcommand{\tr}{\operatorname{tr}}
\newcommand{\ot}{\otimes}
\newcommand{\E}{\mathop{\bf E\/}}
\newcommand{\bx}{\mathbf{x}}
\newcommand{\by}{\mathbf{y}}
\newcommand{\bu}{\mathbf{u}}
\newcommand{\bv}{\mathbf{v}}

% Approximation relation (paper uses \approx_\eps and \simeq_\zeta)
\newcommand{\approxeps}{\approx_{\varepsilon}}
\newcommand{\simgeq}{\simeq_{\zeta}}

% Polynomial spaces (paper uses \polyfunc, \polysub, \polymeas)
\newcommand{\polyfunc}[3]{\operatorname{Poly}(#1,#2,#3)}
\newcommand{\polysub}[3]{\operatorname{PolySub}(#1,#2,#3)}
\newcommand{\polymeas}[3]{\operatorname{PolyMeas}(#1,#2,#3)}


\title{Issue 1099: The Line-169 Replacement Step in the Main-Induction Projectivization Chain}
\author{}
\date{2026-05-04}

\begin{document}
\maketitle

\section{The cited step}

This note concerns the projectivization step in the proof of the main induction
theorem in \cite{Ji2020LowIndividualDegree}.\footnote{Tracked in
\ghissue{1099}. In the TeX source consulted here, the relevant passage is in
\path{references/ldt-paper/inductive_step.tex}, lines~150--180, and the
specific replacement step is lines~167--173.}
At this point in the argument, the paper has already produced polynomial-outcome
measurements $G^{\mathrm A}$ and $G^{\mathrm B}$ with consistency error
$\zeta_1$, projective submeasurements $P^{\mathrm A}$ and
$P^{\mathrm B}$ from orthonormalization, and completed projective measurements
$Q^{\mathrm A}$ and $Q^{\mathrm B}$.

The paper first states the pre-projective consistency estimate
\[
  G^{\mathrm A}_g \otimes I \simeq_{\zeta_1} I \otimes G^{\mathrm B}_g,
  \tag{\path{eq:G-self-consistency}}
\]
then the completion estimate
\[
  G^{\mathrm A}_g \otimes I \approx_{\zeta_2} Q^{\mathrm A}_g \otimes I,
  \tag{\path{eq:G-with-Q-A}}
\]
and finally says that Proposition~\path{prop:triangle-sub}, applied to these two
relations, implies
\[
  Q^{\mathrm A}_g \otimes I \simeq_{\zeta_1} I \otimes G^{\mathrm B}_g.
  \tag{paper line 169}
\]
The same claim is then used, after data processing, in the point-evaluation
transport.

\section{Why the printed inference is incorrect}

The substitution proposition \path{prop:triangle-sub} does not preserve the
original consistency error when one replaces a measurement by an
$\approx_{\epsilon}$-close one.  Its checked content is that from
\[
  A \simeq_{\delta} C
  \qquad\text{and}\qquad
  A \approx_{\epsilon} B
\]
one obtains
\[
  B \simeq_{\delta + \sqrt{\epsilon}} C.
\]
Thus applying it directly to \path{eq:G-self-consistency} and
\path{eq:G-with-Q-A} yields only
\[
  Q^{\mathrm A}_g \otimes I
  \simeq_{\zeta_1 + \sqrt{\zeta_2}}
  I \otimes G^{\mathrm B}_g,
\]
not the printed exact $\zeta_1$ bound.

This is not a formalization artefact.  It is the quantitative content of the
cited proposition itself.  The paper's displayed use of \path{prop:triangle-sub}
therefore drops a square-root term.

\section{A more local corrected estimate}

There is, however, a sharper local repair than the naive
$\zeta_1 + \sqrt{\zeta_2}$ route.  One should compare with the
orthonormalized submeasurement \emph{before} completion.

Orthogonalization gives
\[
  G^{\mathrm A}_g \otimes I
  \approx_{100\zeta_1^{1/4}}
  P^{\mathrm A}_g \otimes I.
\]
For a fixed partner measurement $G^{\mathrm B}$, Proposition
\path{prop:easy-approx-from-approx-delta} bounds the change in diagonal match
mass by the square root of the $\approx$-error.  Specializing that proposition
to the single-question setting gives
\[
  \bigl|
    \sum_g \langle \psi,
      (G^{\mathrm A}_g \otimes G^{\mathrm B}_g)\psi\rangle
    -
    \sum_g \langle \psi,
      (P^{\mathrm A}_g \otimes G^{\mathrm B}_g)\psi\rangle
  \bigr|
  \leq \sqrt{100\zeta_1^{1/4}}
  = 10\zeta_1^{1/8}.
\]
Since \path{eq:G-self-consistency} says that the first match mass is at least
$1-\zeta_1$, one obtains
\[
  \sum_g \langle \psi,
    (P^{\mathrm A}_g \otimes G^{\mathrm B}_g)\psi\rangle
  \geq 1 - \zeta_1 - 10\zeta_1^{1/8}.
\]

Now complete $P^{\mathrm A}$ at a distinguished outcome to obtain
$Q^{\mathrm A}$.  The completion adds the positive residual
$I - \sum_g P^{\mathrm A}_g$ to one outcome, so the diagonal match mass against
the fixed partner measurement $G^{\mathrm B}$ can only increase.  Therefore
\[
  \sum_g \langle \psi,
    (Q^{\mathrm A}_g \otimes G^{\mathrm B}_g)\psi\rangle
  \geq 1 - \zeta_1 - 10\zeta_1^{1/8},
\]
which is equivalent to the corrected consistency estimate
\[
  Q^{\mathrm A}_g \otimes I
  \simeq_{\zeta_1 + 10\zeta_1^{1/8}}
  I \otimes G^{\mathrm B}_g.
\]
The Bob-side mirror is identical after exchanging the tensor factors.

This repaired bound is strictly local to the projectivization step.  It uses
only the already-derived $G$--$P$ orthonormalization comparison and the
positivity of canonical completion.  Quantitatively it is much better than the
direct completed-measurement substitution route
$\zeta_1 + \sqrt{\zeta_2}$, because it inserts the square root before the
completion scalar has been expanded.

\section{Relation with the blueprint and Lean}

The blueprint remark on the Lean line-169 interface now records the active
transport routes and the auxiliary completion fact they use:\footnote{See
\path{blueprint/src/chapter/ch04_projective.tex}, the remark
\path{rem:lean-line169-projectivization-match-mass}.}
the generic $\zeta_1 + \sqrt{\zeta_2}$ substitution route, the sharper local
repaired route $\zeta_1 + 10\zeta_1^{1/8}$ described above, and the checked
completion helper
\texttt{ProjectivizationMatchMassMonotonicity.completeAtOutcomeProj\_left\_matchMass\_ge}
showing that canonical completion adds no further diagonal match-mass loss.

The formal development proves the repaired route by checked declarations.\footnote{The
generic pre-completion transport theorems are
\leanid{MIPStarRE.LDT.MakingMeasurementsProjective.ProjectivizationLine169Repair.leftConsistency\_of\_completion\_and\_sdd}
and
\leanid{MIPStarRE.LDT.MakingMeasurementsProjective.ProjectivizationLine169Repair.rightConsistency\_of\_completion\_and\_sdd}
in
\doclink{MIPStarRE/LDT/MakingMeasurementsProjective/ProjectivizationChain.lean}.
Specializing the $\approx$-error to
\leanid{MIPStarRE.LDT.MakingMeasurementsProjective.orthonormalizationError}
gives
\leanid{MIPStarRE.LDT.MakingMeasurementsProjective.ProjectivizationLine169Repair.leftConsistency\_with\_orthonormalization\_loss}
and its Bob-side mirror.}
These declarations do not prove the paper's exact $\zeta_1$ line.  They prove
the local corrected error $\zeta_1 + 10\zeta_1^{1/8}$.

The generic corrected route $\zeta_1 + \sqrt{\zeta_2}$ remains available as the
literal consequence of \path{prop:triangle-sub}, but propagating that loss all
the way through the final point-transport triangle would force a larger final
envelope at the $1/65536$ scale.  The present formal base-case assembly instead
uses the sharper local pre-completion repair, so the public final envelope keeps
the older $1/40000$ scale.

The former exact-$\zeta_1$ formal interface has now been retired from the active
Lean development.  The formal development instead cleanly distinguishes two
questions:

\begin{itemize}
\item the paper's printed use of \path{prop:triangle-sub}, which is incorrect as
  stated and has the checked local repair above; and
\item the stronger exact-$\zeta_1$ replacement step, which would require an
  additional construction-specific monotonicity theorem not supplied by the
  cited proof and is no longer represented by a live Lean interface.
\end{itemize}

\section{Conclusion}

The verdict is that the paper's line-169 replacement step contains a genuine
quantitative error.  Applying \path{prop:triangle-sub} to the completed
measurement $Q^{\mathrm A}$ does not yield the printed exact $\zeta_1$ bound.

The most local corrected estimate presently available is
\[
  Q^{\mathrm A}_g \otimes I
  \simeq_{\zeta_1 + 10\zeta_1^{1/8}}
  I \otimes G^{\mathrm B}_g,
\]
with the analogous Bob-side statement.  This repair is already formalized and
is substantially sharper than the naive $\zeta_1 + \sqrt{\zeta_2}$ route, but
it is still weaker than the paper's exact line.  The formal public theorem keeps
the older $1/40000$ envelope by using this sharper local repair rather than the
generic corrected route.  Whether the exact $\zeta_1$ transport can be recovered
from a more refined construction-level argument remains a separate open question.

\bibliographystyle{alpha}
\bibliography{references}

\end{document}
